\documentclass[11pt]{article}
\usepackage[margin=1in]{geometry}
\usepackage{listings}
\usepackage{xcolor}
\usepackage{hyperref}
\usepackage{booktabs}
\usepackage[T1]{fontenc}

\definecolor{codegray}{rgb}{0.95,0.95,0.95}
\lstset{
  basicstyle=\ttfamily\small,
  backgroundcolor=\color{codegray},
  breaklines=true,
  frame=single,
  language=,
  showstringspaces=false
}

\newcommand{\TARGET}{\textsf{TARGET}}

\title{Knowledge-Cryptography Trapdoor Extraction over Native QLever Indexes\\
\large Operational Specification (Spec-Proposed; Empirical Tier Labeled \TARGET{}---Unmeasured)}
\author{ggen / kc-trapdoor work order}
\date{2026-10-08}

\begin{document}
\maketitle

\section{Scope and Epistemic Status}

This spec is grounded in exactly one real artifact:
\texttt{/Users/sac/qlever-fortune5-test}, index \texttt{fortune5},
builder \texttt{git-hash 1075455}, \textbf{203 triples}, 47 subjects,
140 objects, 22+7 predicates, all six index permutations on disk
(\texttt{.spo/.pso/.pos/.osp/.sop} and internal \texttt{.pso/.pos}
variants), built
2026-08-26. Everything else---buffer sizes, timeouts, latency---is
spec-proposed and labeled \TARGET{} (unmeasured). No benchmark number in
this document was produced by execution against this corpus. The
billion-triple regime is future work; nothing here claims it.

\section{The Needle-vs-Haystack Invariant}

\subsection{Statement}
Let $G$ be the realization graph (SBB = single-binding-bound, ABB =
all-bindings-bound realizations of the SHACL surface). The invariant:
\emph{realization of a needle must be parameterized by exact cryptographic
coordinates, never by open pattern matching.} Concretely, a needle is
addressed only by the pair
\[
\bigl(\,H_s = \mathrm{BLAKE3}_K(\mathrm{canonical}(s))\,,\;
       P_\sigma = \mathrm{BLAKE3}_K(\mathrm{salted}(p))\,\bigr)
\]
where $K$ is a key held outside the query plane and the hash is computed
off-band. Any query whose WHERE clause leaves a transitive-reachability
variable un-anchored (free $src$/$dst$ of a $^{+}$ path) violates the
invariant and is refused at the admission gate (Section~\ref{sec:tarpits}).

\subsection{CONSTRUCT via merge joins over pre-sorted arrays}
QLever's six permutations (\texttt{fortune5.index.spo}, \texttt{.pso},
\texttt{.pos}, \texttt{.osp}, \texttt{.sop}) are pre-sorted ID arrays.
A keyed coordinate lookup is a binary search: $O(\log N)$ per bound term,
then a merge join across the two operand permutations. The operational
CONSTRUCT is \texttt{queries/needle-realization.construct.rq}: it binds
the needle and salt IRIs from off-band BLAKE3 literals, resolves
$\texttt{kc:salt} \rightarrow \texttt{sh:path}$ against the real index
relations (\texttt{f5:serviceId}, \texttt{f5:region},
\texttt{f5:maxReplicationLagMs}, \texttt{f5:lastRotatedAt},
\texttt{f5:maxAgeDays}, \texttt{f5:P95Latency}, \texttt{f5:P99Latency},
\texttt{f5:availabilitySlaTarget}, \texttt{f5:rotateCredential}), then
joins the SHACL shape surface. The \textbf{<50ms end-to-end latency is a
$\TARGET$}: it has not been measured on this corpus.

\section{Complexity-Asymmetry Argument}

Keyed retrieval is binary search over a sorted permutation:
$O(\log N)$, cheap, per-coordinate. The unkeyed alternative---recovering a
needle by subgraph-isomorphism search over the open graph---is
NP-complete in general \cite{gareyjohnson79}, and the practical adversary
cost is dominated by join-order search over un-anchored transitive
closures. The asymmetry is the security property: holders of $K$ pay
$\log N$; non-holders pay isomorphism. That is the trapdoor.

\begin{thebibliography}{1}
\bibitem{gareyjohnson79} M.~R. Garey and D.~S. Johnson,
\emph{Computers and Intractability}. Freeman, 1979.
\end{thebibliography}

\section{Algorithmic Tar Pits and Fuel Budgets}
\label{sec:tarpits}

\subsection{Saturated cyclic cliques ($b > 100$)}
A trap predicate (here \texttt{f5:rotateCredential}, present in the index
in \texttt{sh:path} position) can be laid out cyclically so that the
$b$-clique closure is saturated. Fuel rule: clique enumeration runs only
for a keyed root (\texttt{queries/tarpit-clique-guard.select.rq}) with
$\mathrm{HAVING}(\ldots > 100)$ as the saturation detector; the
un-anchored closure over the same predicate is refused.

\subsection{Recursive path traps}
The only permitted form of a $^{+}$ traversal is \emph{root-anchored with
explicit \texttt{LIMIT}} (\texttt{queries/tarpit-recursive-path-trap.select.rq},
\texttt{LIMIT 1000}). The refused form is recorded in
\texttt{queries/unanchored-traversal-refusal.select.rq}: a free-variable
$(f5:connects)^+$ enumerates the full closure before any join reduction;
it must be rejected at the admission gate before the engine sees it.

\subsection{QLever configuration parameters}
Buffer sizes, cache caps, and timeouts in
\texttt{qlever/Qleverfile.fortune5} are spec-proposed; the only measured
facts it carries forward from the real Qleverfile are \texttt{NAME},
\texttt{INPUT\_FILES}, \texttt{PORT 7020}, \texttt{ACCESS\_TOKEN}, and
\texttt{SETTINGS\_JSON} (with \texttt{ascii-prefixes-only} changed; the
change is flagged inline). Fuel enforcement is two-layer: (1) admission
gate refuses un-anchored closures with
$\mathrm{KC\_FUEL\_UNANCHORED} = 0$; (2) anchored traversals carry an
explicit \texttt{LIMIT 1000} and a server timeout
$\mathrm{KC\_TIMEOUT\_MS} = 50$, both \TARGET{}.

\section{Zero-Oxigraph Guarantee}

The execution plane is native QLever index artifacts only: reads go to
\texttt{fortune5.index.*} and \texttt{fortune5.internal.index.*} via the
QLever server for index \texttt{fortune5} (builder 1075455). No Oxigraph
library, binary, or embedded store appears in the query path; no RDFlib
fallback execution is permitted for these queries. Any ticket that routes
a KC-trapdoor query through a non-QLever engine is out of scope and
refused.

\section{Falsifiers}
\begin{itemize}
  \item A measured p50 $> 50$ms for the anchored CONSTRUCT on this corpus
        falsifies the \TARGET{} label (and would need re-scoping).
  \item An un-anchored $(f5:connects)^+$ query admitted for execution
        falsifies the fuel-budget invariant (Section~\ref{sec:tarpits}).
  \item Any query artifact executing against a non-native store falsifies
        the zero-Oxigraph guarantee (Section 5).
\end{itemize}

\end{document}
